% ==================== EXERCICE 2 : FONCTIONS ====================

\section{Fonctions \hfill \textit{6 points}}

Pour les fonctions suivantes, donner le \textbf{prototype}, la \textbf{définition} et un \textbf{exemple d'appel} :
    \UPSTIquestion Fonction \texttt{impair()} qui reçoit un entier et qui retourne 1 si cet entier est impair, 0 dans le cas contraire.
\espaceReponse{5cm}
    \UPSTIquestion Fonction \texttt{tronque()} qui reçoit un réel positif et qui retourne la partie entière (sans arrondir).
\espaceReponse{5cm}
    \UPSTIquestion Fonction \texttt{gereDistance()} qui reçoit une distance (réel) et un booléen, et qui retourne la distance modifiée. La distance est augmentée de 10 si le booléen reçu vaut 1, diminuée de 5 si le booléen reçu vaut 0.

\espaceReponse{5cm}
